\documentclass[11pt]{article}
\usepackage[margin=1in]{geometry}
\usepackage{amsmath, amssymb, amsthm}
\usepackage{hyperref}

\newtheorem{theorem}{Theorem}
\newtheorem{definition}{Definition}

\title{Stochastic Processes}
\author{math-foundations / concept 31}
\date{}

\begin{document}
\maketitle

\subsection*{First, an example}

Before any abstract definitions, picture a coin-flip walk. Let
$X_1, X_2, \dots$ be independent with $X_i \in \{-1, +1\}$, each value
having probability $1/2$. Define the partial sums
\[
  S_n \;=\; X_1 + X_2 + \cdots + X_n, \qquad S_0 = 0.
\]
\emph{Read aloud:} ``S sub n equals X sub one plus X sub two plus dot dot dot
plus X sub n, with S sub zero equal to zero.''
Then $E[S_n] = 0$ and $\mathrm{Var}(S_n) = n$, so a typical sample path
wanders to a distance of order $\sqrt{n}$ after $n$ steps. Plotting many
such paths, you see trajectories that jump up or down by exactly one at
each tick and slowly diverge as $\sqrt{n}$. This is the simplest
non-trivial stochastic process and the prototype for everything that
follows.

\section{Definitions}

\begin{definition}[Stochastic process]
Let $(\Omega, \mathcal{F}, P)$ be a probability space and $(S, \mathcal{S})$ a
measurable state space. A \emph{stochastic process} indexed by a set $T$
(typically $T = \mathbb{N}$ or $T = [0, \infty)$) is a collection of
random variables
\[
  X = \{ X_t : t \in T \}, \qquad X_t : \Omega \to S \text{ measurable.}
\]
\emph{Read aloud:} ``X equals the collection of X sub t for t in T, where
each X sub t is a measurable map from big Omega to S.''
For each fixed $\omega \in \Omega$, the function $t \mapsto X_t(\omega)$ is
called a \emph{sample path}.
\end{definition}

\begin{definition}[Markov property]
A stochastic process $\{X_n\}_{n \in \mathbb{N}}$ has the
\emph{Markov property} (with respect to its natural filtration) if for every
$n \in \mathbb{N}$ and every measurable set $A \subseteq S$,
\[
  P\bigl(X_{n+1} \in A \,\bigm|\, X_0, X_1, \dots, X_n\bigr)
  \;=\; P\bigl(X_{n+1} \in A \,\bigm|\, X_n\bigr).
\]
\emph{Read aloud:} ``The probability that X sub n plus one lies in A given
X sub zero, X sub one, through X sub n equals the probability that X sub
n plus one lies in A given X sub n alone.''
Equivalently, given the present $X_n$, the future $X_{n+1}$ is independent
of the past $(X_0, \dots, X_{n-1})$.
\end{definition}

\begin{definition}[Martingale]
Let $\{\mathcal{F}_n\}_{n \in \mathbb{N}}$ be a filtration on
$(\Omega, \mathcal{F}, P)$. A process $\{X_n\}$ is a \emph{martingale}
with respect to $\{\mathcal{F}_n\}$ if
\begin{enumerate}
  \item $X_n$ is $\mathcal{F}_n$-measurable for every $n$ (adaptedness);
  \item $E[|X_n|] < \infty$ for every $n$ (integrability);
  \item $E[X_{n+1} \mid \mathcal{F}_n] = X_n$ almost surely (fair-game
    property).
\end{enumerate}
\end{definition}

\section{Simple Random Walk}

\begin{theorem}
The simple random walk $S_n = X_1 + \cdots + X_n$ with iid
$X_i \in \{-1, +1\}$ each with probability $1/2$ satisfies $E[S_n] = 0$
and $\mathrm{Var}(S_n) = n$.
\end{theorem}

\begin{proof}
Each summand $X_i$ takes values $\pm 1$ with equal probability, so
\[
  E[X_i] \;=\; (+1) \cdot \tfrac{1}{2} + (-1) \cdot \tfrac{1}{2} \;=\; 0,
\]
\emph{Read aloud:} ``Expectation of X sub i equals plus one times one-half
plus minus one times one-half, which equals zero.''
and
\[
  E[X_i^2] \;=\; (+1)^2 \cdot \tfrac{1}{2} + (-1)^2 \cdot \tfrac{1}{2} \;=\; 1.
\]
\emph{Read aloud:} ``Expectation of X sub i squared equals plus one squared
times one-half plus minus one squared times one-half, which equals one.''
Hence $\mathrm{Var}(X_i) = E[X_i^2] - (E[X_i])^2 = 1 - 0 = 1$.

By linearity of expectation,
\[
  E[S_n] \;=\; E\!\left[ \sum_{i=1}^{n} X_i \right]
         \;=\; \sum_{i=1}^{n} E[X_i] \;=\; \sum_{i=1}^{n} 0 \;=\; 0.
\]
\emph{Read aloud:} ``Expectation of S sub n equals expectation of the sum
from i equals one to n of X sub i, which equals the sum from i equals one
to n of expectation of X sub i, equal to zero.''

Because the $X_i$ are independent, all cross-covariances vanish:
$\mathrm{Cov}(X_i, X_j) = 0$ for $i \neq j$. The
variance-of-a-sum identity therefore reduces to a sum of variances:
\[
  \mathrm{Var}(S_n)
  \;=\; \mathrm{Var}\!\left( \sum_{i=1}^{n} X_i \right)
  \;=\; \sum_{i=1}^{n} \mathrm{Var}(X_i)
        + 2 \sum_{i < j} \mathrm{Cov}(X_i, X_j)
  \;=\; \sum_{i=1}^{n} 1 + 0
  \;=\; n.
\]
\emph{Read aloud:} ``Variance of S sub n equals variance of the sum of the
X sub i, which equals the sum of variances plus twice the sum over i less
than j of covariances, which collapses to the sum of ones, equal to n.''
This proves $E[S_n] = 0$ and $\mathrm{Var}(S_n) = n$.
\end{proof}

\paragraph{Remark.} The standard deviation of $S_n$ grows like $\sqrt{n}$,
which is the diffusive scaling underlying Brownian motion (concept 32)
and the noise schedule of diffusion-based generative models.

\end{document}
